% Method section -- AAAI-27.
%
% Changes made for the AAAI kit, all mechanical:
%   - Manual section numbers removed ("2. Method" -> "Method"). main_aaai.tex sets
%     \setcounter{secnumdepth}{2}, so LaTeX numbers these itself and the numbers stay
%     correct if a section is inserted ahead of this one. Put the numbers back by hand
%     only if you also set secnumdepth to 0.
%   - "(Dan and Wu 2026)" -> \citep{dan2026oscillatory}; the entry is in references.bib
%     and aaai2027.bst renders it in AAAI's author-year style.
%   - The hand-written "\section*{References}" and its single entry are gone: references
%     belong in main.tex via \bibliography, not in an \input'ed section.
%   - \clearpage removed -- the kit forbids it ("No page breaks of any kind").
%   - \end{document} removed -- this file is \input'ed.
%   - Appendix removed as requested, with its three \begin{table*} symbol tables.
%     The \label{eq:...} marks are kept so the appendix can be re-attached unchanged.
%
% NEEDS, beyond what aaai2027.sty loads: \usepackage{amsmath} (gathered, aligned,
% \operatorname, \boldsymbol, \eqref). It is not on the kit's forbidden list.

\section{Method}
\label{sec:method}

Our model generates a separate mask for each discovered object by grouping image patches with
similar temporal spike patterns. Inspired by S$^2$-Net \citep{dan2026oscillatory}, it converts an
image into patch activations and uses coupled oscillators to regulate when patches spike.
Figure~\ref{fig:architecture} shows the signal flow.

% Architecture figure -- AAAI-27, spans both columns.
%
% NEEDS in the preamble, beyond the kit and booktabs/amsmath:
%     \usepackage{tikz}
%     \usetikzlibrary{arrows.meta,positioning,fit,backgrounds,calc}
% tikz is not on the kit's forbidden list, and it draws vectors rather than
% including an image file, so the kit's image-format and trim/clip rules do not apply.
%
% WHAT THIS DRAWS, and what it deliberately does not:
%   top-down  theta -> gate g(t) -> dendrite/membrane     (macro modulates micro)
%   bottom-up s(t-1) -> dendritic branch input            (spikes modulate the next branch step)
% There is NO arrow from the spiking layer into the Kuramoto layer, because the trained
% configuration has none: graph_feedback_strength = 0 and spike_pulse_gain = 0
% (hyperparameter.py:95,112), and the rollout raises if either is non-zero
% (sw240/rollout.py:13-15). The graph is built once from gamma and held for the whole
% rollout. Drawing a loop into the oscillators would show something the model does not do.
%
% Symbols match method.tex exactly: gamma, A, phi, theta, g(t), x(t), h, V, s(t), rho, Q.

\begin{figure*}[!t]
\centering
\begin{tikzpicture}[
  font=\small,
  box/.style   = {draw, rounded corners=2pt, align=center, inner sep=4pt, minimum height=8mm},
  stage/.style = {box, fill=black!4, text width=23mm},
  macro/.style = {box, fill=violet!10, draw=violet!60, text width=34mm},
  micro/.style = {box, fill=teal!10,   draw=teal!60,   text width=34mm},
  panel/.style = {box, fill=black!2, text width=31mm},
  flow/.style  = {-{Latex[length=2mm]}, thick},
  td/.style    = {-{Latex[length=2.6mm]}, very thick, violet!70},
  bu/.style    = {-{Latex[length=2.6mm]}, very thick, teal!70},
  lbl/.style   = {font=\scriptsize, align=center, inner sep=1pt}
]

% ---------------------------------------------------------------- left: encoding
\node[stage] (img)  {Image\\$128\times128$};
\node[stage, right=5mm of img] (enc) {Encoder\\8 filters};
\node[stage, right=5mm of enc] (gam) {Patch activations\\$\gamma$, $64\times64$};
\draw[flow] (img) -- (enc);
\draw[flow] (enc) -- (gam);

% ---------------------------------------------------------------- what gamma produces
\node[box, fill=black!4, above right=4mm and 7mm of gam, text width=30mm] (graph)
  {Coupling graph $A$\\[-1pt]{\scriptsize feature $+$ geodesic, top-512}};
\node[box, fill=black!4, below right=4mm and 7mm of gam, text width=30mm] (drive)
  {Sensory drive $\phi$\\[-1pt]{\scriptsize $\pi\tanh(W_\phi\bar{\boldsymbol u})$}};
\draw[flow] (gam.east) -- ++(3mm,0) |- (graph.west);
\draw[flow] (gam.east) -- ++(3mm,0) |- (drive.west);

% ---------------------------------------------------------------- centre: the two scales
\node[macro, right=10mm of graph] (kur)
  {\textbf{Macro-scale}\\[1pt]Kuramoto coordination\\[-1pt]
   {\scriptsize $\theta(t{+}1)=\theta+\Delta t(\omega+f\phi+\kappa\sin(\phi-\theta)+\tfrac{K}{N}\sum A\sin(\theta_j-\theta_i-\alpha))$}};
\node[micro, below=20mm of kur] (snn)
  {\textbf{Micro-scale}\\[1pt]Dendrites $\to$ membrane $\to$ spikes\\[-1pt]
   {\scriptsize $h\to V\to s(t)=g(t)\,\mathrm{ind}(V>\vartheta)$}};

\draw[flow] (graph.east) -- (kur.west);
\draw[flow] (drive.east) -| ([xshift=-6mm]kur.south west) -- ([xshift=-6mm]kur.south west) -- (kur.south west);

% top-down: gating
\draw[td] ([xshift=-9mm]kur.south) -- node[lbl, left=1pt, text=violet!80]
  {\textbf{Top-down}\\gating $g(t)$\\[-1pt]2-step delay\\$x=g\sin\theta$} ([xshift=-9mm]snn.north);

% bottom-up: spike feedback into the next branch input -- stays inside the micro scale,
% which is why it curves back into the same block and not up into the oscillators.
\draw[bu] (snn.east) .. controls +(12mm,6mm) and +(12mm,-6mm) ..
  node[lbl, right=1pt, text=teal!80] {\textbf{Bottom-up}\\spike feedback\\$s(t{-}1)\!\to\!$ branch} (snn.north east);

% ---------------------------------------------------------------- right: readout
\node[panel, right=22mm of kur] (sync)
  {\textbf{Phase synchrony}\\[-1pt]{\scriptsize $\theta$ clusters by object}};
\node[panel, right=22mm of snn] (corr)
  {\textbf{Spike correlation}\\[-1pt]{\scriptsize $\rho\!\to\!Q=\prod_a\max(\rho_a,0)$}};
\node[box, fill=black!4, below=7mm of corr, text width=31mm] (mask)
  {\textbf{Object masks}\\[-1pt]{\scriptsize connected components $\to$ greedy merge}};

\draw[flow, violet!70] (kur.east) -- (sync.west);
\draw[flow, teal!70] (snn.east) -- (corr.west);
\draw[flow] (corr) -- (mask);

% synchrony is a diagnostic, not the scored output
\node[lbl, below=1mm of sync, text=black!55] {diagnostic only};

\end{tikzpicture}
\caption{%
Signal flow. An image is encoded into patch activations $\gamma$, which fix both the coupling
graph $A$ and the sensory drive $\phi$ once per image. Oscillator phases $\theta$ evolve under
Kuramoto coupling on $A$ (macro-scale) and gate the spiking layer through $g(t)$, a two-step
delayed function of the mean patch phase (top-down). Within the spiking layer the previous spike
$s(t{-}1)$ re-enters the dendritic branch input (bottom-up). Object masks are read from
\emph{spike} correlations: $\rho$ over the settled window, combined across the four components
into the affinity $Q$, then connected components and a greedy merge. Phase synchrony is reported
as a diagnostic and is never the scored output.}
\label{fig:architecture}
\end{figure*}


\subsection{From Images to Neural Activations}

\subsubsection{Patch activation.}
We encode an RGB image with eight convolutional filters, normalize and bound the feature maps, and
pool them into patch activations. The patch activation, $\gamma_i$, represents patch $i$. The patch
indices, $i$ and $j$, identify image patches; the component index, $a$, identifies an oscillator
component; and the time index, $t$, identifies a simulation step.

\subsubsection{Patch connection graph.}
To connect related patches, we combine feature agreement with the paths between them.
\begin{equation}
S_{ij}=\frac{\boldsymbol z_i^\top\boldsymbol z_j}{\tau}-r\,d(i,j).
\end{equation}
The connection score, $S_{ij}$, compares the normalized patch embeddings, $\boldsymbol z_i$ and
$\boldsymbol z_j$, using the temperature, $\tau$, and penalizes the geodesic distance, $d(i,j)$,
using the spatial cost weight, $r$. $\tau$ and $r$ are derived from trainable scalars.

We normalize the strongest connections and make their weights reciprocal.
\begin{equation}
\begin{gathered}
 B_{ij}=G\,\operatorname{softmax}_{j\in\mathcal N_i}(S_{ij}),\\
 A_{ij}=\frac{B_{ij}+B_{ji}}{2}.
\end{gathered}
\end{equation}
For each patch $i$, the neighbor set, $\mathcal N_i$, contains the 512 patches with the highest
connection scores, $S_{ij}$. The connection gain, $G$, is derived from a trainable scalar. The
normalized connection weight, $B_{ij}$, is zero for unselected neighbors. The symmetric connection
weight, $A_{ij}$, defines the patch graph.

\subsubsection{Input phase.}
A bounded phase input supplies image information to each oscillator.
\begin{equation}
\begin{gathered}
\boldsymbol u_i=W_{\mathrm{drive}}\gamma_i,\\
\boldsymbol\phi_i=\pi\tanh(W_\phi\,\bar{\boldsymbol u}_i),\qquad
\boldsymbol\theta_i(0)=\boldsymbol\phi_i.
\end{gathered}
\end{equation}
The projection matrix, $W_{\mathrm{drive}}\in\mathbb R^{4\times8}$, and the phase gain,
$W_\phi\in\mathbb R$, are trainable weights. The projected input, $\boldsymbol u_i$, is
standardized to obtain $\bar{\boldsymbol u}_i$. The sensory drive, $\boldsymbol\phi_i$,
initializes the oscillator phase, $\boldsymbol\theta_i(0)$.

\subsection{Recurrent Core}

\subsubsection{Kuramoto layer.}
Oscillators coordinate their phases through intrinsic drift, patch interactions, and sensory drive.
\begin{equation}
\begin{aligned}
\boldsymbol\theta_i(t+1)={}&\boldsymbol\theta_i(t)+\Delta t\Bigl(
\boldsymbol\omega_i+f\boldsymbol\phi_i\\
&+\boldsymbol\kappa_i\odot\sin(\boldsymbol\phi_i-\boldsymbol\theta_i(t))\\
&+\frac{K}{N}\sum_j\bar A_{ij}\sin\bigl(
\boldsymbol\theta_j(t)-\boldsymbol\theta_i(t)-\alpha_{ij}\bigr)\Bigr).
\end{aligned}
\end{equation}
The \textbf{natural frequency}, $\boldsymbol\omega_i$, gives each oscillator its intrinsic rate of
phase change. The \textbf{coupling term}, the sum over $j$, adjusts the phase according to other
patches, using the corrected connection weight, $\bar A_{ij}$, and the phase lag, $\alpha_{ij}$.
The \textbf{drive term} uses the sensory drive, $\boldsymbol\phi_i$, to adjust frequency through
$f$ and to pull the phase toward the input when the drive strength, $\boldsymbol\kappa_i$, is
positive. The coupling scale, $K$, is normalized by the patch count, $N$. The time step, $\Delta
t$, scales the update, and $\odot$ denotes elementwise multiplication. The parameters
$\boldsymbol\omega_i$, $f$, and $\boldsymbol\kappa_i$ are trainable.

\subsubsection{Sinusoidal gating.}
The patch rhythm controls how oscillatory signals enter the spiking core.
\begin{equation}
\begin{gathered}
\bar\theta_i(t)=\frac{1}{4}\sum_{a=1}^4\theta_{ia}(t),\\
g_i(t)=\frac{1+\sin(\bar\theta_i(\max(0,t-2)))}{2}.
\end{gathered}
\end{equation}
\begin{equation}
\boldsymbol x_i(t)=g_i(t)\sin(\boldsymbol\theta_i(t)).
\end{equation}
The mean patch phase, $\bar\theta_i(t)$, summarizes the four oscillator phases in a patch. The
delayed gate, $g_i(t)$, controls signal transmission using the patch phase from two steps earlier.
The gated phase feature, $\boldsymbol x_i(t)$, carries the modulated oscillations to the dendritic
layer.

\subsubsection{Dendritic layer.}
Four branches process the rhythmic input and previous spike before combining their responses.
\begin{equation}
\boldsymbol v_{ia}(t)=U\begin{pmatrix}x_{ia}(t)\\s_{ia}(t-1)\end{pmatrix}+\boldsymbol b.
\end{equation}
The branch input, $\boldsymbol v_{ia}(t)$, combines the gated feature, $x_{ia}(t)$, and the
previous emitted spike, $s_{ia}(t-1)$. The previous spike provides top-down feedback from the
spiking output to the dendritic branches. The projection weight, $U\in\mathbb R^{4\times2}$, and
the bias, $\boldsymbol b\in\mathbb R^4$, are trainable.

The provisional branch state, $\widehat{\boldsymbol h}_{ia}(t)$, balances the previous state and
the current branch input through the retention vector, $\boldsymbol\beta_i$:
\begin{equation}\label{eq:branch-provisional}
\begin{aligned}
\widehat{\boldsymbol h}_{ia}(t)={}&\boldsymbol\beta_i\odot\boldsymbol h_{ia}(t-1)\\
&+(1-\boldsymbol\beta_i)\odot\boldsymbol v_{ia}(t).
\end{aligned}
\end{equation}
\begin{equation}
\begin{aligned}
\boldsymbol h_{ia}(t)={}&\boldsymbol h_{ia}(t-1)
+\bigl(1-\eta_a(1-g_i(t))\bigr)\\
&\quad{}\cdot\bigl(\widehat{\boldsymbol h}_{ia}(t)-\boldsymbol h_{ia}(t-1)\bigr),\\
H_{ia}(t)={}&\sum_{b=1}^{4}h_{iab}(t).
\end{aligned}
\end{equation}
The branch state, $\boldsymbol h_{ia}(t)$, moves toward the provisional branch state,
$\widehat{\boldsymbol h}_{ia}(t)$, according to the gate and the trainable modulation weight,
$\eta_a$. The branch index, $b$, selects an entry of $\boldsymbol h_{ia}(t)$; the membrane input,
$H_{ia}(t)$, sums all four entries.

\subsubsection{Membrane layer.}
The provisional membrane voltage, $\widehat V_{ia}(t)$, balances the previous voltage and the
current membrane input through the retention coefficient, $\lambda_i$, then subtracts a reset
proportional to the previous spike:
\begin{equation}\label{eq:membrane-provisional}
\begin{aligned}
\widehat V_{ia}(t)={}&\lambda_iV_{ia}(t-1)+(1-\lambda_i)H_{ia}(t)\\
&-\vartheta_a s_{ia}(t-1).
\end{aligned}
\end{equation}
\begin{equation}
\begin{aligned}
V_{ia}(t)={}&V_{ia}(t-1)+\bigl(1-\xi_a(1-g_i(t))\bigr)\\
&\quad{}\cdot\bigl(\widehat V_{ia}(t)-V_{ia}(t-1)\bigr).
\end{aligned}
\end{equation}
\begin{equation}
\begin{gathered}
e_{ia}(t)=\operatorname{ind}(V_{ia}(t)>\vartheta_a),\\
s_{ia}(t)=g_i(t)e_{ia}(t).
\end{gathered}
\end{equation}
The membrane voltage, $V_{ia}(t)$, moves toward the provisional voltage, $\widehat V_{ia}(t)$,
according to the gate and the trainable modulation weight, $\xi_a$. The firing threshold,
$\vartheta_a$, is derived from a trainable parameter. The indicator function, $\operatorname{ind}$,
produces the binary firing event, $e_{ia}(t)$, which the gate scales to produce the emitted spike,
$s_{ia}(t)$.

\subsection{Spike Grouping and Mask Generation}

We mean-center and L2-normalize the final 512 emitted spike samples for each oscillator component.
The component correlation, $\rho_{ija}$, sums the products of corresponding samples in the
normalized traces:
\begin{equation}\label{eq:spike-correlation}
\rho_{ija}=\operatorname{clip}\left(
  \sum_{m=1}^{512} r_{ia,m}\,r_{ja,m},-1,1\right),
\end{equation}
where the sample index, $m$, runs over the retained time window. To obtain the patch affinity,
$Q_{ij}$, we replace negative correlations with zero and multiply the resulting values across the
four components. We form connected components by linking patch pairs with $Q_{ij}\geq0.5$, discard
components smaller than 8 patches, a floor selected on CLEVR validation, and treat the largest
remaining component as background. Every patch enters this step as a candidate; no firing-rate
test is applied, so the foreground is defined entirely by the synchrony structure.

We start the retained foreground patches as singleton groups and merge spatially adjacent groups
using
\begin{equation}
\begin{gathered}
E_{ij}=\left(1+\frac{N\bar A_{ij}}{G}\right)\log\frac{Q'_{ij}}{1-Q'_{ij}},\\
\Delta(\mathcal C,\mathcal D)=\sum_{i\in\mathcal C}\sum_{j\in\mathcal D}E_{ij}.
\end{gathered}
\end{equation}
The stabilized affinity, $Q'_{ij}$, determines the sign of the merge evidence, $E_{ij}$. The merge
score, $\Delta(\mathcal C,\mathcal D)$, sums the evidence between the candidate patch groups,
$\mathcal C$ and $\mathcal D$. Groups are merged while their score is positive. Each final group
produces an object mask.
